\chapter{材料把电子交出的能量送往哪里}
\label{chap:feqo-material-response}

电子损失与探测到的光不是同一个量。第~\ref{chap:feqo-transition-kernel} 章已从同一电流和 Green 响应写出二者；现在用平面界面把轨迹积分真正算完。这个例子会出现强电子损失而没有远场光，从而具体说明“激发模式”与“收集到光子”之间缺少哪一步。

\section{可完整积分的平面界面损失谱}
\label{chap:feqo-eels}

令真空占据 \(x>0\)，均匀各向同性材料占据 \(x<0\)。电子以速度 \(v\hat{\vb{z}}\) 沿 \(x=b>0,y=0\) 直线运动。材料的相对介电函数为 \(\varepsilon(\omega)\)。先研究主导波数 \(\kappa\gg\omega/c\) 的非迟滞近场：电场可写成 \(-\nabla\varphi\)，但仍保留材料随频率变化的 \(\varepsilon\)。这个近似要求主要贡献来自 \(b^{-1}\) 或 \(\omega/v\) 所选波数远大于 \(\omega/c\)；若两者接近，必须使用完整 Maxwell Green 张量。

把频域电势沿界面 Fourier 展开为 \(\varphi(x,y,z,\omega)=\int(\dd k_y\dd k_z/(2\pi)^2)e^{\ii(k_yy+k_zz)}\widetilde\varphi(x,\vb{k}_\parallel,\omega)\)。\cref{eq:feqo-electron-current} 对应的电荷密度为 \(\rho=q\delta(x-b)\delta(y)e^{\ii\omega z/v}/v\)。故源的 \(z\) 向 Fourier 变换含 \(2\pi\delta(k_z-\omega/v)\)：电子轨迹直接固定 \(k_z=\omega/v\)。令 \(\kappa=\sqrt{k_y^2+k_z^2}\)，真空半空间的 Poisson 方程为
\[
(\partial_x^2-\kappa^2)\widetilde\varphi=-\widetilde\rho/\varepsilon_0.
\]
在 \(x=0\) 处要求电势连续、法向位移 \(\varepsilon\partial_x\varphi\) 连续，解的散射部分为
\begin{equation}
\widetilde\varphi_{\rm ind}(x,k_y,k_z,\omega)
=\frac{q}{2\varepsilon_0v\kappa}
\frac{1-\varepsilon(\omega)}{1+\varepsilon(\omega)}
e^{-\kappa(x+b)}\,2\pi\delta(k_z-\omega/v),\quad x>0.
\label{eq:feqo-planar-image-potential}
\end{equation}
分母 \(2\kappa\) 来自一维方程的点源导数跳跃；分式是电势的静电反射系数，记为 \(r_{\rm es}\)。若 \(\varepsilon=1\)，材料与真空无区别，诱导势为零。

反射系数值得从边界条件亲手复核。在源与界面之间，源本身的 Fourier 势与它激发的反射势分别正比于 \(e^{\kappa x}\) 与 \(r_{\rm es}e^{-\kappa x}\)；材料一侧的衰减解正比于 \(t_{\rm es}e^{\kappa x}\)。在 \(x=0\) 连续电势得到 \(1+r_{\rm es}=t_{\rm es}\)。令界面无自由表面电荷，法向位移连续给 \(1-r_{\rm es}=\varepsilon t_{\rm es}\)。消去 \(t_{\rm es}\) 才得到 \(r_{\rm es}=(1-\varepsilon)/(1+\varepsilon)\)。在第二个条件中把导数的负号漏掉，会把表面模的条件 \(1+\varepsilon=0\) 错写成别的式子。

由 \(E_z=-\partial_z\varphi_{\rm ind}\) 得 \(E_z=-\ii(\omega/v)\varphi_{\rm ind}\)。把它放入\cref{eq:feqo-eels-work}，对电子经过的长度 \(L\) 积分，再除以 \(L\hbar\omega\)，得到单位轨迹长度、单位角频率的单量子损失概率
\begin{equation}
\frac{\dd^2\Gamma_{\rm EELS}}{\dd z\,\dd\omega}
=\frac{q^2}{2\pi\varepsilon_0\hbar v^2}
\int_{-\infty}^{\infty}\frac{\dd k_y}{2\pi}
\frac{e^{-2b\sqrt{k_y^2+(\omega/v)^2}}}{\sqrt{k_y^2+(\omega/v)^2}}
\bigl[-\operatorname{Im}r_{\rm es}(\omega)\bigr].
\label{eq:feqo-planar-eels}
\end{equation}
对被动材料 \(\operatorname{Im}\varepsilon\geq0\)，直接算分式的虚部给 \(-\operatorname{Im}r_{\rm es}=2\operatorname{Im}\varepsilon/|1+\varepsilon|^2\geq0\)，所以本例的概率非负。这里使用的是由两个边界条件求出的具体散射核；它的正性不能从任意散射 Green 张量单独正定来猜。

还可将横向波数积分做到底。令 \(k_y=(\omega/v)\sinh u\)，则 \(\dd k_y/\sqrt{k_y^2+(\omega/v)^2}=\dd u\)。定义修正 Bessel 函数 \(K_0(x)=\int_0^\infty e^{-x\cosh u}\dd u\)，于是
\begin{equation}
\frac{\dd^2\Gamma_{\rm EELS}}{\dd z\,\dd\omega}
=\frac{q^2}{2\pi^2\varepsilon_0\hbar v^2}
K_0\!\left(\frac{2\omega b}{v}\right)
\bigl[-\operatorname{Im}r_{\rm es}(\omega)\bigr].
\label{eq:feqo-planar-eels-k0}
\end{equation}
左边单位为 \(\si{\second\per\metre}\)，与右边的 \(q^2/(\varepsilon_0\hbar v^2)\) 相同；乘 \(\dd z\dd\omega\) 后才是无量纲概率。\(K_0\) 随 \(\omega b/v\) 增大而衰减，表示离界面越远，高频近场越难被电子采到。

衰减还可以定量估计。对 \(x=2\omega b/v\gg1\)，积分 \(K_0(x)=\int_0^\infty e^{-x\cosh u}\dd u\) 主要来自 \(u\simeq0\)。在这里 \(\cosh u=1+u^2/2+O(u^4)\)，所以把积分的主贡献换成高斯积分给
\[
K_0(x)=\sqrt{\frac{\pi}{2x}}e^{-x}
\left[1+O(x^{-1})\right].
\]
因此把电子离界面距离增加 \(\Delta b\)，高频信号的主要指数因子减少 \(e^{-2\omega\Delta b/v}\)。这正是实验中高频近场对轨迹位置格外敏感的原因；若束斑宽度与 \(v/(2\omega)\) 可比，就应对不同 \(b\) 的概率平均，而不能只代入平均轨迹。

若取 Drude 材料 \(\varepsilon(\omega)=1-\omega_p^2/[\omega(\omega+\ii\gamma)]\)，弱阻尼 \(\gamma\ll\omega_p\) 时 \(1+\varepsilon\) 在 \(\omega\simeq\omega_p/\sqrt2\) 附近变小，损失谱出现表面模峰；\(\gamma\) 控制峰宽。无限平面沿 \(z\) 保持平移对称，电子给出的 \(k_z=\omega/v>\omega/c\) 不在真空传播光锥内，所以在此理想几何中远场阴极发光为零。材料吸收或表面波仍可让\cref{eq:feqo-planar-eels-k0} 为正。这是损失与发光不能混同的一个完整反例。

这条 Drude 谱的形状可以直接从\cref{eq:feqo-planar-eels-k0} 数值积出来。\cref{fig:planar-drude-eels} 取 \(\hbar\omega_p=9.0\,\si{\electronvolt}\)、\(\hbar\gamma=0.08\,\si{\electronvolt}\)、\(v=0.70c\)、\(b=20\,\si{\nano\metre}\)，实线是含推迟效应的完整 Weyl 积分，虚线标出无损极限下的同步极点。谱峰落在 \(5.25\,\si{\electronvolt}\)，几乎正对虚线给出的 \(5.23\,\si{\electronvolt}\)；阻尼把本该发散的极点削成有限高度 \(2.6\times10^{-4}\,\si{\per\electronvolt\per\nano\metre}\)，峰形明显偏向高能侧——\(6.0\,\si{\electronvolt}\) 处仍有峰值的四分之一，而低能侧 \(4.5\,\si{\electronvolt}\) 已不到峰值的百分之一。虚线与实线的重合说明摆位预测由 \(1+\varepsilon=0\) 给出，峰高和峰宽才需要 \(b\) 与 \(\gamma\) 入场。

\begin{figure}[htbp]
\centering
\includegraphics[width=0.86\linewidth]{planar_drude_eels.pdf}
\caption{真空侧掠射电子对 Drude 半空间的单位长度 EELS 谱。实线为完整 Weyl 积分数值结果，虚线为无损同步极点 \(\hbar\omega=\hbar\omega_p/\sqrt{1-\varepsilon_{\rm sync}}\simeq5.23\,\si{\electronvolt}\)；参数取 \(\hbar\omega_p=9.0\,\si{\electronvolt}\)、\(\hbar\gamma=0.08\,\si{\electronvolt}\)、\(\beta=0.70\)、\(b=20\,\si{\nano\metre}\)。纵轴单位为 \(\si{\per\electronvolt\per\nano\metre}\)。}
\label{fig:planar-drude-eels}
\end{figure}

峰宽的说法可由复介电函数而非图像推出。写 \(\omega_s=\omega_p/\sqrt2\)、\(\delta=\omega-\omega_s\)，并取 \(|\delta|,\gamma\ll\omega_s\)。逐项展开得
\[
1+\operatorname{Re}\varepsilon\simeq\frac{4\delta}{\omega_s},
\qquad
\operatorname{Im}\varepsilon\simeq\frac{2\gamma}{\omega_s},
\qquad
-\operatorname{Im}r_{\rm es}\simeq
\frac{\gamma\omega_s}{4\delta^2+\gamma^2}.
\]
在 \(K_0(2\omega b/v)\) 相对这个窄峰变化缓慢时，损失峰的全半高宽近似为 \(\gamma\)。若 \(b\) 很大使 \(K_0\) 在峰宽内显著变化，测得的峰形会被轨迹选择重新加权；若迟滞不可忽略，静电反射系数本身也必须换成 Maxwell 边界问题的结果。

\begin{exercise}
从界面两侧电势连续和 \(\varepsilon\partial_x\varphi\) 连续推导\cref{eq:feqo-planar-image-potential} 中的反射系数，并求无损 Drude 极限的峰频率。
\end{exercise}


若样品有边缘、缺陷或光栅，界面上的非辐射场可以向远场散射。判定阴极发光时必须继续算能量穿过哪个面、进入哪个角度和偏振，而不能把材料 Green 函数的全部虚部直接称为光子数。

\section{远场光子是总损失中的哪一部分}
\label{chap:feqo-cl}

平面界面的电子可激发很强的近场，却因平移动量过大而没有远场光。若样品有限、边缘破坏平移对称性，或另有光栅提供动量，式~\eqref{eq:feqo-cl-collection} 才可能给非零计数。实际计算先用同一电流源求 \(\mathbf E=\ii\mu_0\omega\mathbf G\mathbf j_e\)，再取其远场系数 \(\mathbf F(\hat{\mathbf n},\omega)\)。这个顺序保留了材料响应与观测投影的区别。

把这个操作写成核的形式，就能看见探测器究竟投影了什么。对远处的观测点定义 \(\mathbf G_\infty(\hat{\mathbf n},\mathbf r';\omega)=\lim_{r\to\infty}r e^{-\ii kr}\mathbf G(r\hat{\mathbf n},\mathbf r';\omega)\)，前提是该方向存在通常的 \(e^{\ii kr}/r\) 远场展开。则
\[
\mathbf F(\hat{\mathbf n},\omega)
=\ii\mu_0\omega\int\dd^3r'\,
\mathbf G_\infty(\hat{\mathbf n},\mathbf r';\omega)
\mathbf j_e(\mathbf r',\omega).
\]
同一电子电流在式~\eqref{eq:feqo-trajectory-green-kernel} 中与 \(\operatorname{Im}\mathbf G\) 作二次型，给所有损失通道；在这里先由 \(\mathbf G_\infty\) 投影到指定方向，再取模平方，只给能到远场的部分。吸收模与导波模没有该方向的远场系数，因而不会凭空出现在 CL 计数中。

对包住样品的闭合曲面，时域 Poynting 定理把电子交出的能量分成介质吸收、向外辐射和最终仍储存在局域场中的三项。若脉冲前后局域储能回到原值，频率分辨的能量收支可写为 \(W_{\rm loss}=W_{\rm abs}+W_{\rm rad}\)，但只有在同一物理区域、同一入射条件且模式间交叉项正确处理后才可逐频使用。定义 \(\eta_{\rm rad}=W_{\rm rad}/(W_{\rm rad}+W_{\rm abs})\) 与 \(\eta_{\rm coll}=W_{\rm coll}/W_{\rm rad}\)，二者分别回答“多少能量成为远场光”和“多少远场光进入本探测器”。光学数值孔径只能改变第二项，不能恢复已变成热的能量。

由式~\eqref{eq:feqo-cl-collection}，改变收集角不一定只改变总强度：若两个模式有不同角分布，探测器对它们的相对权重也会改变。比较电子损失谱与发光谱时，必须给出收集角、偏振与效率；把一个未标定的发光峰直接解释成材料态密度会混入仪器选择。

一个小颗粒的偶极近似可把角度权重算完。假设颗粒远小于波长，电子在颗粒位置产生的频域驱动场为 \(\mathbf E_e(\omega)\)，颗粒的线性极化率为体积量纲的复数 \(\alpha(\omega)\)。若响应只沿 \(z\) 方向，诱导偶极矩为 \(p_z(\omega)=\varepsilon_0\alpha(\omega)E_{e,z}(\omega)\)。真空偶极远场满足
\[
\mathbf F(\hat{\mathbf n},\omega)
=\frac{k^2}{4\pi\varepsilon_0}
\bigl[\hat{\mathbf n}\times(\mathbf p\times\hat{\mathbf n})\bigr],
\qquad k=\omega/c.
\]
把 \(\mathbf p=p_z\hat{\mathbf z}\) 代入，\(|\mathbf F|^2=k^4|p_z|^2\sin^2\theta/(16\pi^2\varepsilon_0^2)\)。式~\eqref{eq:feqo-cl-collection} 因而预言角分布为 \(\sin^2\theta\)，而不是各方向相同。若理想探测器收集以 \(+z\) 为轴、半角 \(\theta_{\max}\) 的圆锥，且效率在圆锥内为一，则因 \(\dd\Omega=2\pi\sin\theta\dd\theta\)，收集到的辐射比例为
\begin{equation}
\eta_{\rm cone}=
\frac{\int_0^{\theta_{\max}}\sin^3\theta\dd\theta}
{\int_0^\pi\sin^3\theta\dd\theta}
=\frac{2-3\cos\theta_{\max}+\cos^3\theta_{\max}}4.
\label{eq:feqo-dipole-collection-fraction}
\end{equation}
半球收集给 \(1/2\)，很窄的轴向圆锥因偶极轴上恰有辐射零点而远小于它的立体角比例。这个算例从已知电子驱动场 \(E_{e,z}\) 经材料极化率算到远场角分布与收集计数；它依赖偶极、线性响应和真空远场条件，不能直接用在尺寸可比波长的颗粒或复杂衬底上。

\begin{exercise}
设一个模式把总能量的 \(80\%\) 送到远场，探测器收集其中 \(25\%\)。求 \(\eta_{\rm rad}\eta_{\rm coll}\)，并说明改变数值孔径能否改变模式内部的吸收功率。
\end{exercise}


空间周期性又提供了额外动量。它能打开前述平面界面中缺失的辐射通道，也使电子与混合模式在特定速度和角度下强耦合。

\section{周期结构如何补上缺少的动量}
\label{chap:feqo-structured-media}

令结构沿电子轨迹方向的周期为 \(a\)，其倒格矢为 \(G_m=2\pi m/a\)。场的一个 Bloch 分量写作 \(E_m(z,\omega)=u_m e^{\ii(k_z+G_m)z}\)。它与运动电子在\cref{eq:feqo-beta} 中的相位相乘，并在长度 \(L\) 的对称作用区积分，得到
\begin{equation}
\int_{-L/2}^{L/2}\dd z\,
e^{\ii(k_z+G_m-\omega/v)z}
=L\sinc\!\left[\frac{(k_z+G_m-\omega/v)L}{2}\right].
\label{eq:feqo-periodic-phase-matching}
\end{equation}
长结构使相位匹配窗口收窄到 \(|k_z+G_m-\omega/v|\lesssim1/L\)。光栅给或取整数个 \(\hbar G_m\) 后，原先落在光锥外的电子近场可能耦合到传播光子；但 \(u_m=0\) 时积分仍为零，说明相位匹配只给允许条件，强度还取决于模式场与电流的重叠。

为什么非要周期结构不可，看一张几何图就清楚。\cref{fig:light-cone-synchronism} 把光锥 \(\omega/c=|k_z|\) 与电子线 \(\omega/c=\beta_0|k_z|\) 画在同一个 \((k_z,\omega)\) 平面上。取 \(\beta_0=0.70\)，两条直线只在原点相遇，其余各处中间的阴影区就是 \(\omega/v=k_z>\omega/c\) 的部分：电子一路能提供的纵向波数全都落在光锥以外，因此在平移对称的理想几何里它无法辐射。要让 \(k_z+G_m\) 落进光锥，只能在横轴上平移整数个倒格矢，这就是\cref{eq:feqo-periodic-phase-matching} 里必须有非零 \(u_m\) 的几何含义。

\begin{figure}[htbp]
\centering
\includegraphics[width=0.72\linewidth]{light_cone_synchronism.pdf}
\caption{光锥与电子线的相对位置，取 \(\beta_0=0.70\)。横轴为 \(k_z/k_0\)，纵轴为 \(\omega/(ck_0)\)；实线是光锥 \(\omega/c=|k_z|\)，虚线是电子线 \(\omega/c=\beta_0|k_z|\)，阴影区是电子能提供的波数全部落在光锥之外的区域，只能靠倒格矢平移进入。}
\label{fig:light-cone-synchronism}
\end{figure}

例如出射光在真空中与电子轨迹夹角为 \(\theta\)，其纵向波数为 \(k_z=(\omega/c)\cos\theta\)。把它代入相位匹配式，取 \(m>0\) 与 \(\lambda=2\pi c/\omega\)，得
\begin{equation}
\lambda=\frac{a}{m}
\left(\frac{c}{v}-\cos\theta\right).
\label{eq:feqo-smith-purcell-dispersion}
\end{equation}
这条关系说明电子速度、光栅周期和收集角共同决定可见波长。由于 \(v<c\)，括号在任何实角都为正；改变 \(a\) 会按比例移动谱线，而改变探测角会在同一结构上选不同波长。有限作用长度时并不存在严格的单一角度：\cref{eq:feqo-periodic-phase-matching} 的第一零点满足 \(|\Delta k|=2\pi/L\)，在固定 \(\omega\) 下对应 \(|\Delta\cos\theta|\simeq\lambda/L\)。这个估计只给相位匹配宽度，实际角分布还乘模式外耦合、颗粒尺寸及收集效率。若无相应的 \(u_m\)，即使满足\cref{eq:feqo-smith-purcell-dispersion} 也没有这一级辐射。

材料中的两个裸模也可相互混合。先在无损、只含两个归一化模的模型中定义
\[
H_{\rm mix}=\hbar\begin{pmatrix}\omega_a&g\\g^*&\omega_b\end{pmatrix}.
\]
解二阶特征方程得 \(\omega_\pm=(\omega_a+\omega_b)/2\pm\sqrt{(\omega_a-\omega_b)^2/4+|g|^2}\)。当 \(\omega_a=\omega_b\) 时分裂为 \(2|g|\)；本征向量决定电子电流对上下两支各有多大重叠。给裸模另加衰减率 \(\kappa_a,\kappa_b\) 后，有效矩阵的对角元可写成 \(\omega_{a,b}-\ii\kappa_{a,b}/2\)，但它不是封闭系统的 厄米 哈密顿量；复本征值描述峰位和线宽的近似，而峰是否能被看成两条，还取决于耦合、线宽、探测窗口和背景干涉。

上面的 \(\omega_\pm\) 在无量纲坐标里就是一对双曲线。\cref{fig:hopfield-avoided-crossing} 的横轴是失谐 \((\bar\omega_c-\omega_m)/|g_{\rm cm}|\)，纵轴是 \((\omega-\bar\omega)/|g_{\rm cm}|\)，两条灰线是未混合的重整化类光子模与裸物质模，在原点交叉；两条实线是混合后的两支，在失谐为零处不相交，而是一上一下停在 \(\pm1\)，即最小劈裂正好等于 \(2|g_{\rm cm}|\)。失谐推到 \(\pm8\) 时实线几乎贴在虚线上，说明远离共振时两个裸模各自复位；可分辨的劈裂只存在于 \(|\bar\omega_c-\omega_m|\) 与 \(|g_{\rm cm}|\) 同量级的窗口内。

\begin{figure}[htbp]
\centering
\includegraphics[width=0.76\linewidth]{hopfield_avoided_crossing.pdf}
\caption{RWA 下的反交叉。横轴为重整化失谐 \((\bar\omega_c-\omega_m)/|g_{\rm cm}|\)，纵轴为相对频率 \((\omega-\bar\omega)/|g_{\rm cm}|\)；虚线为重整化类光子模，点线为裸物质模，两条实线为混合后的本征频率 \(\omega_\pm\)，在 \(\bar\omega_c=\omega_m\) 处的最小劈裂为 \(2|g_{\rm cm}|\)。}
\label{fig:hopfield-avoided-crossing}
\end{figure}

把运动学与收集通道放在同一实验里，电子可激发满足\cref{eq:feqo-periodic-phase-matching} 的内部模；材料损耗决定其中多少变热；边缘或光栅耦合决定多少成为辐射；最后光学系统只读收集到的角度与偏振。这几个比例分别变化，所以改变电子速度、结构周期或数值孔径不能被解释为同一个参数的效应。

\begin{exercise}
在 \(k_z=0\) 的一维光栅里解 \(\omega/v=G_m\) 对周期 \(a\) 的条件，并说明有限长度 \(L\) 将该条件展宽多少。
\end{exercise}

